\begin{alist}
\item Από το σχήμα είναι φανερό ότι η γραφική της παράσταση προκύπτει από επανάληψη του τμήματος της που αντιστοιχεί στο διάστημα $[0, \pi]$, οπότε η περίοδος της $f$ είναι $T = \pi$.
Επιπλέον οι τεταγμένες των σημείων της γραφικής της παράστασης περιέχονται στο διάστημα $[-3, 3]$, οπότε η ελάχιστη τιμή της είναι ίση με $-3$ και η μέγιστη είναι ίση με 3.

\item Η συνάρτηση είναι της μορφής $f(x) = \rho\eta\mu(\alpha x)$ με $\alpha, \rho > 0$, οπότε έχει περίοδο $T = \dfrac{2\pi}{\alpha}$ ελάχιστη τιμή $-\rho$ και μέγιστη ίση με $\rho$. Έτσι, έχουμε $\dfrac{2\pi}{\alpha} = \pi$, οπότε $\alpha = 2$ και $\rho = 3$.

\item Είναι: $g(x) = x^4 - 2x^2 + 5 = (x^2 - 1)^2 + 4 \ge 4$ και η ισότητα $g(x) = 4$ ισχύει όταν $x^2 = 1$ δηλαδή όταν $x = -1$ ή $x = 1$. Άρα ο αριθμός 4 είναι η ελάχιστη τιμή της $g$.

\item Επειδή για κάθε $x \in \mathbb{R}$ ισχύει $f(x) \le 3$ και $g(x) \ge 4$, η εξίσωση $f(x) = g(x)$ είναι αδύνατη. Άρα οι γραφικές παραστάσεις των $f, g$ δεν έχουν κοινό σημείο.
\end{alist}
